% ===========================================================================
%  固件 EKF 组合导航理论与算法分析设计报告 —— 公共导言区
% ===========================================================================
\documentclass[UTF8,a4paper,12pt]{ctexart}

% ---------- 页面与段落 ----------
\usepackage[top=2.6cm,bottom=2.6cm,left=2.7cm,right=2.5cm]{geometry}
\usepackage{setspace}
\onehalfspacing
\setlength{\parindent}{2em}
\setlength{\parskip}{0.2em}
\setlength{\emergencystretch}{2em}
\raggedbottom
\clubpenalty=10000
\widowpenalty=10000
\displaywidowpenalty=10000
\usepackage{microtype}

% ---------- 页眉页脚与标题层级 ----------
\usepackage{fancyhdr}
\usepackage{titlesec}
\usepackage{caption}
\setcounter{tocdepth}{2}
\setcounter{secnumdepth}{3}
\setlength{\headheight}{15pt}
\pagestyle{fancy}
\fancyhf{}
\fancyhead[L]{\small\kaishu 纵列双发矢量推力飞行器固件 EKF 组合导航}
\fancyhead[R]{\small\kaishu 理论与算法分析设计报告}
\fancyfoot[C]{\small\thepage}
\renewcommand{\headrulewidth}{0.35pt}
\renewcommand{\footrulewidth}{0pt}
\fancypagestyle{plain}{%
  \fancyhf{}
  \fancyfoot[C]{\small\thepage}
  \renewcommand{\headrulewidth}{0pt}
  \renewcommand{\footrulewidth}{0pt}
}
\titleformat{\section}{\Large\bfseries}{\thesection}{0.8em}{}
\titleformat{\subsection}{\large\bfseries}{\thesubsection}{0.8em}{}
\titleformat{\subsubsection}{\normalsize\bfseries}{\thesubsubsection}{0.8em}{}
\titlespacing*{\section}{0pt}{2.0ex plus .5ex minus .2ex}{1.0ex}
\titlespacing*{\subsection}{0pt}{1.5ex plus .4ex minus .2ex}{0.7ex}
\titlespacing*{\subsubsection}{0pt}{1.2ex plus .3ex minus .2ex}{0.5ex}

% ---------- 数学与定理 ----------
\usepackage{amsmath,amssymb,amsthm}
\usepackage{mathtools}
\usepackage{bm}
\usepackage{bbm}
\theoremstyle{definition}
\newtheorem{definition}{定义}[section]
\theoremstyle{plain}
\newtheorem{theorem}{定理}[section]
\theoremstyle{remark}
\newtheorem{remark}{备注}[section]
\newtheorem*{notation}{符号约定}

% ---------- 图表 ----------
\usepackage{graphicx}
\usepackage{float}
\usepackage{placeins}
\usepackage{tikz}
\usepackage{pgfplots}
\pgfplotsset{compat=1.18}
\usepgfplotslibrary{groupplots}
\usetikzlibrary{arrows.meta,shapes,positioning,calc,patterns,angles,quotes,3d,backgrounds,fit}
\usepackage{booktabs}
\usepackage{array}
\usepackage{tabularx}
\usepackage{multirow}
\captionsetup{font=small,labelfont=bf,labelsep=space,justification=centering}

% ---------- 代码 ----------
\usepackage{listings}
\lstset{
  basicstyle=\ttfamily\small,
  frame=single,
  numbers=left,
  numberstyle=\tiny\color{gray},
  breaklines=true,
  columns=flexible,
  keepspaces=true,
  language=C++,
  commentstyle=\color{teal},
  keywordstyle=\color{blue},
}

% ---------- 引用与链接 ----------
\usepackage[numbers,sort&compress]{natbib}
\bibliographystyle{gbt7714-numeric}
\usepackage[hidelinks,unicode]{hyperref}
\usepackage{cleveref}
\hypersetup{
  pdftitle={纵列双发矢量推力飞行器固件 EKF 组合导航：理论与算法分析设计报告},
  pdfauthor={LShang},
  pdfsubject={15 状态误差状态 EKF 组合导航的理论推导、算法设计与仿真验证},
  pdfkeywords={EKF, 捷联惯导, GNSS 组合导航, 双子样圆锥补偿, 静止辅助, 误差状态}
}
\crefname{equation}{式}{式}
\crefname{figure}{图}{图}
\crefname{table}{表}{表}
\crefname{section}{节}{节}
\crefname{subsection}{节}{节}
\crefname{subsubsection}{节}{节}
\crefname{appendix}{附录}{附录}
\renewcommand{\refname}{参考文献}

% ---------- 列表 ----------
\usepackage{enumitem}
\setlist{nosep,leftmargin=2.2em}

% ---------- 自定义命令 ----------
% 常用矩阵/矢量记号
\newcommand{\vp}[1]{\boldsymbol{#1}}          % 粗体矢量
\newcommand{\m}[1]{\boldsymbol{#1}}            % 粗体矩阵
\newcommand{\q}[1]{\mathfrak{#1}}              % 四元数
% 注意：\dtheta / \dphi 本身已是粗体符号；\sk{} 直接括起已加粗参数，不再二次加粗
\newcommand{\dtheta}{\boldsymbol{\theta}}
\newcommand{\dphi}{\boldsymbol{\phi}}
\newcommand{\dx}{\delta\vp{x}}
\newcommand{\Cbn}{\m{C}_b^n}
\newcommand{\Cnb}{\m{C}_n^b}
\newcommand{\I}[1]{\m{I}_{#1}}
\newcommand{\E}[1]{\mathbb{E}\left[#1\right]}
\newcommand{\diag}{\operatorname{diag}}
\newcommand{\sk}[1]{[\,#1\,\times]}
\newcommand{\tr}{\operatorname{tr}}
\newcommand{\T}{^{\top}}
\newcommand{\inv}{^{-1}}
\newcommand{\Exp}{\operatorname{Exp}}
\newcommand{\Log}{\operatorname{Log}}
\newcommand{\atan}[1]{\operatorname{atan}(#1)}
\newcommand{\atanh}[1]{\operatorname{atanh}(#1)}
\newcommand{\atanhtwo}{\operatorname{atan2}}
\newcommand{\cross}{\times}
\newcommand{\code}[1]{\texttt{#1}}

% 文件引用
\newcommand{\file}[1]{\textcolor{blue}{\texttt{#1}}}

% 表格列宽工具
\newcolumntype{L}[1]{>{\raggedright\arraybackslash}p{#1}}
\newcolumntype{C}[1]{>{\centering\arraybackslash}p{#1}}
\newcolumntype{R}[1]{>{\raggedleft\arraybackslash}p{#1}}

\usepackage{makecell}
\renewcommand{\theadfont}{\small\bfseries}
